% !TEX root = ../thesis.tex
\chapter{Kết quả thực nghiệm}

Chương này báo cáo kết quả thực nghiệm định lượng chi tiết của mô hình HKL-ViT5 và hai mô hình đối chứng trên tập dữ liệu đánh giá 300 mẫu báo chí tiếng Việt độc lập.
Mục~4.1 trình bày thiết lập đánh giá và cấu trúc định dạng dữ liệu đầu ra;
Mục~4.2 báo cáo kết quả nhóm Lexical Overlap (ROUGE);
Mục~4.3 báo cáo kết quả nhóm Semantic Quality (BERTScore);
Mục~4.4 báo cáo kết quả nhóm Factual Consistency (Tính nhất quán sự thật);
Mục~4.5 báo cáo kết quả nhóm Performance (Hiệu năng suy luận); và Mục~4.6 tổng kết chương.

\section{Thiết lập đánh giá và định dạng dữ liệu}

\subsection{Tập dữ liệu kiểm thử và các mô hình so sánh}

Quá trình đánh giá thực nghiệm được vận hành trên tập kiểm thử độc lập gồm 300 mẫu bài báo thu thập từ trang báo điện tử \textit{tienphong.vn}. Ba mô hình được đưa vào đánh giá đối đầu trực tiếp trên cùng 300 mẫu kiểm thử này:

\begin{enumerate}[label=\arabic*.]
  \item \textbf{HKL-ViT5 (Phương pháp đề xuất):} Mô hình ViT5-base cải tiến kết hợp cơ chế Selective Long-Context, Keyword Conditioning và Heuristic Factuality Reranking.
  \item \textbf{\texttt{Nishikyen/vit5-vietnamese-news}:} Mô hình ViT5-base fine-tune chuẩn trên dữ liệu tin tức tiếng Việt (độ dài ngữ cảnh 256/512 tokens).
  \item \textbf{\texttt{thnhan/sft\_model}:} Mô hình ViT5-base tinh chỉnh theo phương pháp SFT.
\end{enumerate}

\subsection{Cấu trúc định dạng đầu ra dữ liệu}

Kết quả sinh ra của mỗi mô hình trên từng văn bản kiểm thử được lưu trữ dưới dạng các tệp văn bản mã hóa UTF-8 với định dạng tiêu chuẩn:

\begin{tcolorbox}[colback=gray!5!white,colframe=gray!75!black,title=Cấu trúc định dạng tệp kết quả tóm tắt đầu ra (.txt)]
\small
\begin{verbatim}
INDEX:
<index_bài_báo>
GUID:
<mã_định_danh_guid>
TITLE:
<tiêu_đề_bài_báo>
VĂN BẢN GỐC:
<toàn_văn_nội_dung_bài_báo_gốc>
TÓM TẮT DO MODEL SINH:
<nội_dung_bản_tóm_tắt_sinh_ra_bởi_mô_hình>
THỜI GIAN TÓM TẮT:
<thời_gian_xử_lý_tính_bằng_giây>
---Hết---
\end{verbatim}
\end{tcolorbox}

\section{Kết quả nhóm 1: Lexical Overlap (ROUGE)}

Bảng~\ref{tab:results_rouge} báo cáo điểm ROUGE-1, ROUGE-2 và ROUGE-L (tính theo F1 Macro trung bình trên 300 mẫu kiểm thử chung).

\begin{table}[H]
  \centering
  \caption{Kết quả đánh giá nhóm Lexical Overlap (ROUGE F1 Macro trên 300 mẫu). \textbf{In đậm}: Kết quả tốt nhất.}
  \label{tab:results_rouge}
  \begin{tabularx}{\linewidth}{X CCC}
    \toprule
    Mô hình & ROUGE-1 & ROUGE-2 & ROUGE-L \\
    \midrule
    \textbf{HKL-ViT5 (Ours)}               & \textbf{0.5184} & \textbf{0.2932} & \textbf{0.3434} \\
    \texttt{Nishikyen/\allowbreak vit5-vietnamese-news} & 0.4707 & 0.2510 & 0.3148 \\
    \texttt{thnhan/sft\_model}             & 0.4321 & 0.2227 & 0.2773 \\
    \bottomrule
  \end{tabularx}
\end{table}

\noindent \textbf{Phân tích kết quả:}
Mô hình HKL-ViT5 dẫn đầu vượt trội trên cả 3 chỉ số ROUGE. Cụ thể, HKL-ViT5 đạt ROUGE-1 0.5184 (cao hơn \texttt{Nishikyen} +4.77\% và cao hơn \texttt{thnhan} +8.63\%). Điểm ROUGE-2 đạt 0.2932 (tăng +4.22\%) và ROUGE-L đạt 0.3434 (tăng +2.86\%). Kết quả này minh chứng rằng việc nén ngữ cảnh toàn văn bằng MMR và Keyword Conditioning giúp mô hình khôi phục được nhiều n-gram quan trọng từ bài gốc mà các mô hình bị cắt ngắn đầu vào bỏ sót.

\section{Kết quả nhóm 2: Semantic Quality (BERTScore)}

Bảng~\ref{tab:results_bertscore} báo cáo các chỉ số tương đồng ngữ nghĩa BERTScore (sử dụng mô hình mã hóa ngữ cảnh \texttt{xlm-roberta-base}).

\begin{table}[H]
  \centering
  \caption{Kết quả đánh giá nhóm Semantic Quality (BERTScore với \texttt{xlm-roberta-base}). \textbf{In đậm}: Kết quả tốt nhất.}
  \label{tab:results_bertscore}
  \begin{tabularx}{\linewidth}{X CCC}
    \toprule
    Mô hình & BERTScore Precision & BERTScore Recall & BERTScore F1 \\
    \midrule
    \textbf{HKL-ViT5 (Ours)}               & \textbf{0.8996} & \textbf{0.9020} & \textbf{0.9008} \\
    \texttt{Nishikyen/\allowbreak vit5-vietnamese-news} & 0.8972 & 0.8896 & 0.8934 \\
    \texttt{thnhan/sft\_model}             & 0.8960 & 0.8790 & 0.8874 \\
    \bottomrule
  \end{tabularx}
\end{table}

\noindent \textbf{Phân tích kết quả:}
Ở độ đo ngữ nghĩa BERTScore, HKL-ViT5 đạt F1 0.9008, thiết lập mức hiệu năng cao nhất trong 3 mô hình. Đặc biệt, chỉ số BERTScore Recall đạt 0.9020 (so với 0.8896 của \texttt{Nishikyen} và 0.8790 của \texttt{thnhan}), khẳng định rằng bản tóm tắt do HKL-ViT5 sinh ra bao hàm ngữ nghĩa toàn diện hơn đối với bài viết gốc.

\section{Kết quả nhóm 3: Factual Consistency (Tính nhất quán sự thật)}

Bảng~\ref{tab:results_factuality} báo cáo các chỉ số về tính trung thực thực thể, tính chính xác số liệu, tỷ lệ mâu thuẫn và tỷ lệ ảo giác (sử dụng mô hình NLI \texttt{MoritzLaurer/mDeBERTa-v3-base-mnli-xnli} và trích xuất thực thể \texttt{underthesea}).

\begin{table}[H]
  \centering
  \caption{Kết quả đánh giá nhóm Factual Consistency. Entity Precision, Number Consistency: càng cao càng tốt; Contradiction Rate, Hallucination Rate: càng thấp càng tốt. \textbf{In đậm}: Tốt nhất.}
  \label{tab:results_factuality}
  \begin{tabularx}{\linewidth}{X CCCC}
    \toprule
    Mô hình & Entity Precision (\%) & Number Consistency (\%) & Contradiction Rate (\%) & Hallucination Rate (\%) \\
    \midrule
    \textbf{HKL-ViT5 (Ours)}               & \textbf{81.81\%} & \textbf{98.54\%} & \textbf{2.11\%} & \textbf{39.04\%} \\
    \texttt{Nishikyen/\allowbreak vit5-vietnamese-news} & 79.16\% & 96.46\% & 2.90\% & 45.53\% \\
    \texttt{thnhan/sft\_model}             & 76.57\% & 97.55\% & 3.99\% & 48.45\% \\
    \bottomrule
  \end{tabularx}
\end{table}

Bảng~\ref{tab:results_factuality_supp} bổ sung các chỉ số thống kê về xác suất kéo theo logic và số lượng thực thể/con số được xác nhận thực tế.

\begin{table}[H]
  \centering
  \caption{Thống kê số lượng thực thể, con số và xác suất kéo theo NLI.}
  \label{tab:results_factuality_supp}
  \begin{tabularx}{\linewidth}{X CCC}
    \toprule
    Chỉ số & HKL-ViT5 & Nishikyen & thnhan \\
    \midrule
    Thực thể được xác nhận / Tổng sinh & \textbf{706 / 863} & 581 / 734 & 513 / 670 \\
    Con số được xác nhận / Tổng sinh & \textbf{472 / 479} & 409 / 424 & 438 / 449 \\
    Xác suất Entailment trung bình ($P_{\text{entail}}$) & \textbf{0.6063} & 0.5416 & 0.5153 \\
    \bottomrule
  \end{tabularx}
\end{table}

\noindent \textbf{Phân tích kết quả:}
HKL-ViT5 thể hiện ưu thế vượt trội trong việc bảo tồn tính xác thực:
\begin{itemize}
  \item Entity Precision đạt 81.81\% (xác nhận 706/863 thực thể chính xác), vượt trội so với \texttt{Nishikyen} (79.16\%) và \texttt{thnhan} (76.57\%).
  \item Number Consistency đạt mức tiệm cận tuyệt đối 98.54\% (472/479 con số chính xác), hạn chế tối đa việc nhầm lẫn dữ liệu báo chí.
  \item Contradiction Rate giảm xuống còn 2.11\% (thấp hơn \texttt{thnhan} gần một nửa).
  \item Hallucination Rate giảm đáng kể xuống 39.04\% (so với 45.53\% của \texttt{Nishikyen} và 48.45\% của \texttt{thnhan}), chứng minh hiệu quả kéo giảm ảo giác của bộ Heuristic Factuality Reranker.
\end{itemize}

\section{Kết quả nhóm 4: Performance (Hiệu năng suy luận)}

Bảng~\ref{tab:results_performance} báo cáo các chỉ số thống kê về thời gian suy luận trung bình (tính bằng giây trên mỗi bài báo).

\begin{table}[H]
  \centering
  \caption{Kết quả thống kê thời gian suy luận (Inference Time). \textbf{In đậm}: Tốc độ nhanh nhất.}
  \label{tab:results_performance}
  \begin{tabularx}{\linewidth}{X CCCC}
    \toprule
    Mô hình & Mean (s) & Median (s) & P95 (s) & Std (s) \\
    \midrule
    \textbf{HKL-ViT5 (Ours)}               & 8.0535 & 5.7128 & 18.5620 & 4.7382 \\
    \texttt{Nishikyen/\allowbreak vit5-vietnamese-news} & \textbf{3.5313} & \textbf{3.5053} & \textbf{4.5211} & \textbf{0.5923} \\
    \texttt{thnhan/sft\_model}             & 6.1358 & 5.9035 & 8.1466 & 1.3003 \\
    \bottomrule
  \end{tabularx}
\end{table}

\noindent \textbf{Phân tích kết quả:}
Mô hình \texttt{Nishikyen} có thời gian suy luận nhanh nhất (trung bình 3.53 giây/bài) do chỉ chạy suy luận Beam Search đơn trên chuỗi ngắn (256/512 tokens). Trong khi đó, HKL-ViT5 mất trung bình 8.05 giây/bài (trung vị 5.71 giây). Sự gia tăng thời gian này là hợp lý vì HKL-ViT5 phải thực hiện quy trình nén ngữ cảnh 1024 tokens qua MMR và sinh $N=5$ ứng viên tóm tắt để qua bộ Reranker.

\section{Tổng kết chương}

Thực nghiệm trên 300 mẫu kiểm thử độc lập đã chứng minh tính hiệu quả toàn diện của HKL-ViT5: chiến thắng tuyệt đối trên cả 3 chỉ số ROUGE, đạt BERTScore F1 cao nhất (0.9008), đạt độ chính xác con số 98.54\% và kéo giảm tỷ lệ ảo giác xuống 39.04\%. Sự gia tăng thời gian suy luận (8.05s) là sự đánh đổi hoàn toàn hợp lý để đổi lấy chất lượng và tính xác thực dữ kiện báo chí.
